% discussion.tex (국문 v11)
\section{고찰}\label{sec:discussion}

\subsection{물리 기준값과 전이성}
결과는 물리 지식이 모델에 들어가는 방식이 ML이 지역 Stefan 모델을 넘어서는지, 그 개선이 새 지역으로 전이되는지를 결정함을 보여 준다. 물리 잔차 구성(식~(\ref{eq:residual}))에서 Stefan 항은 해마다의 융해기 기후에 대한 반응을 맡아 물리적으로 제한된 형태로 새 지역과 새 해에 외삽되고, 학습기는 토양과 지형에 따른 실효 계수의 편차를 보정한다. 이 역할 분담을 뒷받침하는 결과는 네 가지이다. 이득에는 지형과 토양 공변량이 필요하였고, Stefan 오차가 가장 큰 조건에서 가장 컸으며, 잔차 학습기의 기여 분석은 레나 델타에서 지형, 캐나다에서 토양 유기탄소를 가리켰고, Stefan 유사라벨은 학습기를 Stefan 모델 쪽으로 끌어당겨 정확도를 높이지 못하였다. 기후 반응을 라벨에서 직접 배워야 하는 직접 ML은 한 지역의 라벨을 모두 쓸 때 물리 잔차 모델보다 점추정 오차가 낮았으나, 레나 델타를 학습에서 제외하면 오차가 Stefan 모델보다 \ndrefLeregDzabs{} cm 컸다. 이 대비는 과정 지식이 주로 자료가 부족하거나 외삽이 필요할 때 성능을 높인다는 호소학 · 수문학의 물리 기반 학습 연구와 일치한다\cite{Read2019,Jia2021,Willard2022}. 대부분의 영구동토 지역에는 ALT 관측이 적거나 없으므로, 지도 작성에는 관측이 많은 지역에서 라벨을 모두 쓴 정확도보다 대상 지역 라벨 없이 유지되는 정확도가 더 중요하다.

\subsection{선행 연구와의 비교}
기존 ALT 지도 연구는 기후 입력, 검증 설계, 비교 기준에서 이 연구와 다르다. ML 제품은 여러 해에 걸쳐 모은 관측을 다년 예측 변수와 연결하고, 각자의 지점 집합에서 무작위 또는 지점 단위 교차검증 결과를 보고한다\cite{Ran2022a,Liu2024,Wei2026}. 이 연구의 자료에서 무작위 교차검증은 직접 ML의 공간 블록 오차를 \noptimLeDz--\noptimCaDz{} cm 과소 추정하였고, 레나 델타를 학습에서 제외하면 직접 ML의 오차는 \nrmseLeregDz{} cm로 커져 다른 지역 계수를 쓴 Stefan 모델(\nrmseLeregPz{} cm)보다 컸다. 따라서 무작위 또는 지점 단위 분할에서 얻은 오차는 학습 관측이 없는 지역에서 ML 제품의 정확도를 나타내지 않는다. Gautam 등\cite{Gautam2025}은 2014년에 관측한 알래스카 CALM 68개 지점에서 랜덤 포레스트와 Stefan 모델을 비교하였고, 무작위 검정 분할에서 랜덤 포레스트의 결정계수는 학습 0.84에서 검정 0.24로 떨어진 반면 Stefan 모델은 0.54를 유지하였다. 이 연구는 이 비교를 세 지역의 2010--2024년 셀-연도와 지역 홀드아웃으로 확장하고, 학습기가 Stefan 모델을 대체하지 않고 보정하도록 두 접근을 결합하였다. ESA CCI Permafrost 같은 과정 기반 제품은 ALT 관측으로 학습하지 않으며\cite{Westermann2024,Westermann2025}, 외부 CALM 지점에서 ESA CCI v5의 오차(\nftAllCcifivezero{} cm)는 지역 라벨 5개를 쓴 우리 모델의 오차(\nftAllRsfive{} cm)보다 컸다.

모든 모델과 제품을 같은 셀 · 같은 해에서 제품 학습 지점 마스크와 함께 채점하면 이러한 차이를 비교할 수 있다. 제품의 지역별 순위는 제품의 학습 자료를 따랐다. 제품이 인근 관측으로 학습한 지역(알래스카, 캐나다)에서는 제품이 우리 모델만큼 정확하였고, 그렇지 않은 지역(레나 델타)에서는 얇은 활동층을 과대 추정하여 레나 델타 라벨 없이 학습한 모델보다도 덜 정확하였다. 잔차 학습기는 제품을 기준값으로 쓸 때에도 정확도를 높였으므로(식~(\ref{eq:prod})), 지역 라벨과 잔차 학습을 결합하면 물리 모델뿐 아니라 기존 제품도 그 제품이 학습하지 않은 지역에 맞게 보정할 수 있다. 같은 셀 · 같은 해 비교 절차와 학습 지점 마스크는 다른 제품과 지역에도 적용할 수 있다.

\subsection{새 지역 · 새 연도 적용 절차}
라벨 수 실험 결과는 새 지역에 방법을 적용하는 순서로 정리된다. 지역 라벨이 없을 때는 다른 지역 계수를 쓴 Stefan 모델이나 다른 지역으로 학습한 물리 잔차 모델이 직접 ML보다 안전한 출발점이다. 직접 ML의 레나 델타 오차는 대상 지역 라벨 없이 \ntlrmseLeDszero{} cm였다. 레나 델타 라벨 25개는 물리 잔차 모델의 오차를 \ntlrmseLeRszero{} cm에서 \ntlrmseLeRstwfive{} cm로 줄여 원천 계수 Stefan 모델(\ntlrmseLePzzero{} cm)보다 낮추었고, 외부 CALM 지역에서는 라벨 5개로 오차가 \nftAllRszero{} cm에서 \nftAllRsfive{} cm로 감소하였다. 이는 수문학의 지역화 연구에서 보고된 소수 지역 관측의 가치와 일치한다\cite{Pool2019}. 필요한 라벨 수는 새 지역이 학습 범위에서 얼마나 떨어져 있는지에 따라 달라진다. 융해 깊이 범위가 넓은 캐나다에서는 라벨 25개로 재보정하면 오차가 증가하였으므로, 지역의 토양과 지형 범위를 포괄하는 라벨 25--100개가 적절한 목표이다. 각 지도는 명시된 해의 ERA5-Land 기후로 계산되고 적용 가능 영역과 예측 구간을 함께 제공하므로, 같은 워크플로로 감시용 연도별 지도를 만들고, 같은 해의 관측이나 제품과 비교하며, 추가 관측이 필요한 적용 가능 영역 밖의 지역을 표시할 수 있다.

\subsection{한계와 향후 과제}
알래스카 라벨의 절반이 2014년 자료여서 이 지역에서는 관측 연도 입력의 효과를 판정하지 못하였다. 레나 델타 라벨은 한 자료원의 \nnblkLe{}개 블록에서 나와 신뢰 구간이 넓다. 기후 입력은 월자료이고 토양 공변량은 정적이어서 단기 기상과 토양 수분 변화는 반영되지 않는다. 주로 산지와 영구동토 지대의 남쪽 가장자리에 있는 약 3 m보다 두꺼운 활동층은 학습 범위 밖이며, 외부 시험은 잔차 학습기가 지역 라벨 없이는 이러한 조건으로 전이되지 않음을 보여 준다. 일 단위 기후 입력, 적설 · 토양 수분 정보, CALM과 집중 조사 자료가 함께 있는 서시베리아 같은 네 번째 지역의 라벨이 다음 과제이다.
